\section{Taxonomy of Partisan Fairness Metrics}
\label{sec:definition}

\subsection{Descriptive vs.\ Normative Metrics}

The six metrics in the L-series divide into two categories with distinct behavioral
properties and distinct roles in litigation.

\textbf{Descriptive metrics} measure the partisan properties of a specific map
given observed vote shares. They answer the question: ``What is the partisan
outcome of this map as drawn?'' They do not specify what the outcome \textit{should}
be. Efficiency Gap, Mean-Median Difference, Partisan Bias, and Declination are
all descriptive. A descriptive metric with value zero does not mean the map is
``fair'' in any normative sense---it means the map has a particular property
(vote-symmetry, zero mean-median gap, etc.) that is consistent with a large class
of different underlying geographic configurations.

\textbf{Normative metrics} specify a target outcome. They answer the question:
``By how much does this map deviate from the outcome that a fair process should
produce?'' The Gallagher Index (proportionality) and responsiveness (from the
seats-votes curve) are normative---they require a prior specification of what
``fair'' means (proportional seat allocation or competitive responsiveness)
before deviations can be measured. The normative content makes them more
contestable in court but also more directly relevant to the constitutional
question of what state law requires.

\subsection{Metric Definitions}

\subsubsection{Efficiency Gap}

For a $k$-district congressional plan, let $w_i^D$ be the Democratic wasted votes
in district $i$ (votes above 50\%+1 for Democratic winners; all votes for Democratic losers)
and $w_i^R$ be Republican wasted votes symmetrically. Total votes cast are $T$.
The Efficiency Gap is:
\[
\text{EG} = \frac{\sum_i w_i^D - \sum_i w_i^R}{T}.
\]
$\text{EG} > 0$ indicates Republican advantage (Democratic votes are wasted at
a higher rate). The range is approximately $[-1, 1]$; values outside $[-0.15, 0.15]$
are rare in empirical data. \citet{stephanopoulosMcGhee2015} proposed an 8\% threshold
as a presumptive unconstitutionality standard. SCOTUS in \textit{Gill v.\ Whitford}
did not adopt this threshold.\footnote{\textit{Gill v.\ Whitford}, 585 U.S.\ 48 (2018).}

\subsubsection{Mean-Median Difference}

Let $v_1 \leq v_2 \leq \cdots \leq v_k$ be the sorted Democratic vote shares across
$k$ districts, $\bar{v}_D = k^{-1}\sum_i v_i$ the mean, and $\tilde{v}_D$ the median.
The Mean-Median Difference is:
\[
\text{MM} = \bar{v}_D - \tilde{v}_D.
\]
Positive MM indicates Democratic voters are geographically sorted into fewer,
higher-margin districts (mean exceeds median). Introduced by \citet{wang2016}.
Range approximately $[-0.10, +0.10]$ for congressional delegations.

\subsubsection{Partisan Bias}

Let $S(v)$ denote the Democratic seat share when Democrats receive statewide
two-party vote share $v$, estimated via uniform swing. Partisan Bias is:
\[
\text{Bias} = S(0.50) - 0.50.
\]
Positive Bias favors Democrats (they would win more than half the seats at exactly
half the votes). Negative Bias favors Republicans. Part of the symmetry framework
developed by \citet{grofmanKing2007}. Range approximately $[-0.20, +0.20]$ for
state congressional delegations.

\subsubsection{Declination}

Let $\bar{v}_D^+$ be the mean Democratic vote share in Democratic-won districts
(those with $v_i > 0.50$) and $\bar{v}_R^-$ be the mean Democratic vote share in
Republican-won districts (those with $v_i < 0.50$). Define angles:
\[
\theta_D = \arctan(2\bar{v}_D^+ - 1), \qquad \theta_R = \arctan(1 - 2\bar{v}_R^-).
\]
Declination is:
\[
\delta = \frac{2(\theta_D - \theta_R)}{\pi}.
\]
Positive $\delta$ indicates Republican advantage through cracking (Republican-won
districts are clustered near 50\%, while Democratic wins are pushed to high-margin
territory). Introduced by \citet{warrington2018}.

\subsubsection{Seats-Votes Curve and Responsiveness}

The seats-votes curve $S(v)$ maps statewide Democratic vote share to expected
Democratic seat share. Estimated from observed district vote shares by applying
a uniform swing $\delta = v - \bar{v}_D$ to each district:
\[
S(v) = \frac{1}{k}\sum_{i=1}^k \mathbf{1}[v_i + \delta > 0.50].
\]
Responsiveness is the slope at the 50\% pivot:
\[
R = \left.\frac{dS}{dv}\right|_{v=0.50}.
\]
$R \approx 2.0$ is the competitive ideal (a one-point vote gain translates to a
two-point seat gain near the median). $R < 1.5$ indicates suppressed competition.
The full $S(v)$ curve is the sufficient statistic from which EG and Partisan Bias
are recoverable analytically (see L.5).

\subsubsection{Gallagher Proportionality Index}

For a two-party system with Democratic vote share $v_D$ and seat share $s_D$:
\[
\text{LSq} = \sqrt{\tfrac{1}{2}\left[(v_D - s_D)^2 + (v_R - s_R)^2\right]}.
\]
$\text{LSq} = 0$ indicates perfect proportionality. Higher values indicate
greater disproportionality. Unlike the other metrics, LSq measures deviation
from a specific normative target (proportionality) rather than measuring
a property of the vote distribution itself.

\subsection{Relationships Among Metrics}

The six metrics are mathematically related in ways that constrain their simultaneous
behavior. Three key relationships:

\textbf{EG and Bias from $S(v)$.} As shown in L.5, EG equals the signed area
between the $S(v)$ curve and the diagonal $S(v) = v$, while Partisan Bias equals
$S(0.50) - 0.50$. The full seats-votes curve is therefore the sufficient statistic
from which both scalar metrics can be computed.

\textbf{MM and EG.} Both metrics are positive when Democratic voters are
concentrated in fewer, higher-margin districts. Their ratio depends on the
distribution of district vote shares; for symmetric distributions, $\text{MM}
\approx \text{EG}/2k$ (see \citealt{mcGhee2017}).

\textbf{Declination and EG.} Declination captures the geometric structure that
EG aggregates into a scalar. A map with high positive Declination (cracking of
Democratic voters) will typically have high positive EG, but the converse does
not hold: EG can be elevated by packing alone without producing the angular
geometry that Declination measures.

\subsection{Decision Table for Expert Witnesses}

Table~\ref{tab:decision-table} provides a six-row decision table mapping the
primary legal theory in each type of state constitutional claim to the most
appropriate metric or metrics.

\begin{table}[ht]
\centering
\caption{Expert Witness Metric Selection Guide for Post-\textit{Rucho} State Litigation}
\label{tab:decision-table}
\begin{threeparttable}
\begin{tabular}{@{}lllp{4cm}@{}}
\toprule
Legal Theory & State Example & Primary Metric & Key Limitation \\
\midrule
Vote symmetry / equal treatment & NC, OH & Partisan Bias & Requires swing model \\
Wasted-vote asymmetry & WI, PA & Efficiency Gap & Sensitive to wave elections \\
Geographic sorting vs.\ manipulation & NC, WI & Mean-Median & Sorting confound \\
Packing-cracking geometry & Any & Declination & No court adoption (2026) \\
Competitive responsiveness & General & $R$ from $S(v)$ & Uniform swing assumption \\
Proportionality floor & PA, NC & Gallagher LSq & No federal requirement \\
\bottomrule
\end{tabular}
\begin{tablenotes}
\small
\item Note: Multiple metrics should be reported together to resist metric-shopping criticism.
The seats-votes curve ($S(v)$) subsumes EG and Bias and is the preferred single exhibit.
\end{tablenotes}
\end{threeparttable}
\end{table}
